% ID: SEL-D0006-OA-E03
% ESTADO: cotejada. Segunda pasada visual completa y borrador compilado revisados.
% Fuente: Selectividad/sel_2019_matematicas_aplicadas/Reserva B/Exámenes Matemáticas Aplicadas CCSS - A y B.pdf
% SHA-256: f9236ec66af0dbd3fbd0231fdcaa6d922425c04d48178ea0737073e76e8eb7d5
% Página: 1; opción A; ejercicio 3.

Para tratar cierta enfermedad, en un hospital se utilizan tres fármacos
distintos, A, B y C, administrándose a cada enfermo un solo fármaco. El
30\,\% de los pacientes es tratado con el fármaco A, el 50\,\% es tratado
con el B y el resto con el fármaco C. La probabilidad de que la enfermedad
se cure con el fármaco A es de 0.6, de que se cure con el fármaco B es de
0.8 y de que se cure con el fármaco C es de 0.7. Se elige al azar un paciente
de ese hospital con esa enfermedad.
\begin{enumerate}[label=\alph*)]
  \item \textbf{(1.5 puntos)} Calcule la probabilidad de que el paciente se
  cure.
  \item \textbf{(1 punto)} Sabiendo que el paciente se ha curado,
  \textquestiondown{}cuál es la probabilidad de que haya sido tratado con
  el fármaco A?
\end{enumerate}
